\chapter{Convex analysis and Legendre duality}
\index{convex analysis}\index{Legendre--Fenchel transform}\index{Alexandrov measure}
\label{app:convexity}

This appendix is a reference sheet for the convex operations used throughout
the book.
It fixes three boundaries that are easy to blur.
Legendre--Fenchel duality requires closedness for exact recovery;
Alexandrov Monge--Amp\`ere measure is not multiplied by \(n!\);
and Alexandrov--Fenchel is invoked only for finite-dimensional convex bodies.
The algebraic factor \(n!\) is restored by
\Cref{app:normalizations}.

The first uses occur in \Cref{chap:topological-tools} for energy
duality, \Cref{chap:bj-homogeneity-bdivisors} for test curves, and
\Cref{chap:transcendental-okounkov-dh} for concave transforms and
Duistermaat--Heckman measures.

\section{Closed convex functions}
\label{sec:appD-closed-convex}

Let \(N\) be a finite-dimensional real vector space, let
\(M=N^\vee\), and write the pairing as \(\langle m,x\rangle\).
The extended real line is
\[
 \overline{\mathbf R}=\mathbf R\cup\{-\infty,+\infty\}.
\]
Undefined expressions such as \(+\infty-\infty\) are never formed.

\begin{definition}[Proper closed convex function]
\label{def:appD-proper-closed}
A function \(f:N\to\mathbf R\cup\{+\infty\}\) is proper if it is never
\(-\infty\) and is finite somewhere.
It is closed convex if its epigraph
\[
 \operatorname{epi}(f)
 :=\{(x,r)\in N\times\mathbf R:r\geq f(x)\}
\tag{D.1}\label{eq:appD-epigraph}
\]
is closed and convex.
Its effective domain is
\(\operatorname{dom}(f)=\{f<+\infty\}\).
\end{definition}

For a proper convex function, its lower semicontinuous closure
\(\overline f\) is the function whose epigraph is the closure of
\(\operatorname{epi}(f)\).

\begin{lemma}\label{lem:appD-closure-rules}
The supremum of any nonempty family of convex functions is convex.
The pointwise limit of a decreasing net of convex functions is convex,
provided it never takes the value \(-\infty\).
The maximum of finitely many closed convex functions is closed convex.
\end{lemma}

\begin{proof}
An epigraph of a supremum is the intersection of the epigraphs.
For the decreasing-net assertion, let \(f=\inf_a f_a\).
Given \(x,y\), \(s\in[0,1]\), and \(\varepsilon>0\), directedness supplies
one index \(a\) such that
\[
 f_a(x)<f(x)+\varepsilon,\qquad
 f_a(y)<f(y)+\varepsilon.
\]
Here there is something to prove only when both \(f(x)\) and \(f(y)\) are
finite; otherwise the convexity inequality has right-hand side
\(+\infty\).
Convexity of \(f_a\), followed by \(\varepsilon\downarrow0\), gives convexity
of \(f\).
Finally, the epigraph of a finite maximum is a finite intersection of closed
convex sets.
\end{proof}

\begin{warning}{Infima}
An arbitrary infimum of convex functions need not be convex.
For example, the minimum of
\(x\mapsto(x-1)^2\) and \(x\mapsto(x+1)^2\) has two wells.
Whenever the main text takes a convex rooftop, it means the largest closed
convex minorant of the pointwise infimum.
\end{warning}

\section{Legendre--Fenchel duality}
\label{sec:appD-legendre}

\begin{definition}[Conjugate and subgradient]
\label{def:appD-conjugate}
For \(f:N\to\mathbf R\cup\{+\infty\}\), define
\[
 f^*(m)=\sup_{x\in N}\{\langle m,x\rangle-f(x)\}.
\tag{D.2}\label{eq:appD-conjugate}
\]
If \(f(x)\) is finite, its subdifferential at \(x\) is
\[
 \partial f(x)
 :=\{m\in M:
 f(y)\geq f(x)+\langle m,y-x\rangle\ \text{for every }y\in N\}.
\tag{D.3}\label{eq:appD-subgradient}
\]
\end{definition}

\begin{proposition}\label{prop:appD-fenchel-equality}
If \(f\) is proper, then for every \(x\in N\) and \(m\in M\),
\[
                       f(x)+f^*(m)\geq\langle m,x\rangle.
\tag{D.4}\label{eq:appD-fenchel}
\]
Equality holds if and only if
\(m\in\partial f(x)\), and then also \(x\in\partial f^*(m)\).
\end{proposition}

\begin{proof}
The inequality follows by retaining \(x\) among the competitors in
\eqref{eq:appD-conjugate}.
Equality says
\[
 \langle m,y\rangle-f(y)
 \leq \langle m,x\rangle-f(x)
\]
for every \(y\), which rearranges to \eqref{eq:appD-subgradient}.
The same equality, read with the variables reversed, gives the last claim.
\end{proof}

\begin{theorem}\label{thm:appD-fenchel-moreau}
Let \(f:N\to\mathbf R\cup\{+\infty\}\) be proper and convex.
Then \(f^*\) is closed convex and
\[
                          f^{**}=\overline f.
\tag{D.5}\label{eq:appD-biconjugate}
\]
In particular, \(f^{**}=f\) exactly when \(f\) is lower semicontinuous.
\end{theorem}

\begin{proof}
The function \(f^*\) is a supremum of continuous affine functions.
Fenchel's inequality gives \(f^{**}\leq f\), hence
\(f^{**}\leq\overline f\).
Fix \(x_0\) and \(a<\overline f(x_0)\).
The point \((x_0,a)\) lies strictly below the closed convex epigraph of
\(\overline f\).
The affine-minorant form of strict epigraph separation produces
\(m\in M\) and \(c\in\mathbf R\) such that
\[
 \langle m,x\rangle-c\leq\overline f(x)
 \quad\text{for all }x,\qquad
 \langle m,x_0\rangle-c>a.
\]
For completeness, this form follows from ordinary strong separation as
follows.  Upward closedness forces the vertical coefficient, with the
orientation chosen so that the epigraph lies above the separating
hyperplane, to be nonnegative.  If it is positive, normalize it to one.  If
it is zero, combine the horizontal separator with a positive-vertical
supporting hyperplane at a relative-interior point of the proper epigraph;
a sufficiently small positive multiple preserves strictness at
\((x_0,a)\), and normalization again gives the displayed affine minorant.
The first inequality says \(f^*(m)\leq c\), so \(f^{**}(x_0)>a\).
Let \(a\uparrow\overline f(x_0)\).
This is the finite-dimensional argument behind
\cite[Theorem~12.2]{Roc70}.
\end{proof}

\begin{example}[Indicator and support function]
\label{ex:appD-indicator-support}
For a nonempty closed convex set \(P\subset M\), let
\(\iota_P=0\) on \(P\) and \(+\infty\) off \(P\).
Its conjugate is
\[
 h_P(x)=\sup_{m\in P}\langle m,x\rangle.
\tag{D.6}\label{eq:appD-support}
\]
\Cref{thm:appD-fenchel-moreau} gives
\[
                         h_P^*=\iota_P.
\tag{D.7}\label{eq:appD-support-dual}
\]
Compactness of \(P\) is equivalent to finiteness of \(h_P\) everywhere.
\end{example}

\begin{proposition}\label{prop:appD-legendre-tests}
For proper functions \(f,g\) and \(c\in\mathbf R\),
\[
\begin{aligned}
 f\leq g&\Longrightarrow f^*\geq g^*,\\
 (f+c)^*&=f^*-c,\\
 (f(\,\cdot-x_0))^*(m)&=f^*(m)+\langle m,x_0\rangle.
\end{aligned}
\tag{D.8}\label{eq:appD-legendre-tests}
\]
\end{proposition}

\begin{proof}
All three statements follow directly from
\eqref{eq:appD-conjugate}.
Only finite \(c\) is allowed.
\end{proof}

\section{Concave test curves and rays}
\label{sec:appD-test-curves}

For \(g:\mathbf R\to\mathbf R\cup\{-\infty\}\), closed concavity means that
\(-g\) is closed convex.
Set
\[
 \mathcal Lg(t)=\sup_{\tau\in\mathbf R}\{g(\tau)+t\tau\},
 \qquad t\geq0.
\tag{D.9}\label{eq:appD-concave-transform}
\]

\begin{theorem}\label{thm:appD-concave-duality}
Let \(g\) be proper, upper semicontinuous, concave, and nonincreasing.
Then \(\mathcal Lg\) is convex on \([0,+\infty)\) and
\[
                  g(\tau)=\inf_{t\geq0}
                    \{\mathcal Lg(t)-t\tau\}.
\tag{D.10}\label{eq:appD-concave-inversion}
\]
Conversely, if \(\ell\) is proper closed convex on \([0,+\infty)\), extended by
\(+\infty\) to \(t<0\), then
\[
 g(\tau)=\inf_{t\geq0}\{\ell(t)-t\tau\}
\]
is closed concave and nonincreasing, and \(\mathcal Lg=\ell\).
\end{theorem}

\begin{proof}
Apply \Cref{thm:appD-fenchel-moreau} to \(f=-g\).
Because \(g\) is nonincreasing, all slopes of \(f\) are nonnegative, so
the effective domain of \(f^*\) lies in \([0,+\infty)\).
Now
\[
 f^*(t)=\sup_\tau\{t\tau+g(\tau)\}=\mathcal Lg(t),
\]
and biconjugation gives \eqref{eq:appD-concave-inversion}.
The converse is the same computation in reverse.
\end{proof}

\begin{definition}[Abstract test curve]
\label{def:appD-test-curve}
Let \(\mathcal C\) be an ordered convex cone of extended-real functions on a
set \(X\).
A concave test curve is a family
\(\psi=(\psi_\tau)_{\tau\in\mathbf R}\subset\mathcal C\) satisfying
\[
 \psi_{(1-s)\tau_0+s\tau_1}
 \geq(1-s)\psi_{\tau_0}+s\psi_{\tau_1},
 \qquad 0<s<1,
\tag{D.11}\label{eq:appD-test-curve-concavity}
\]
that is nonincreasing and pointwise upper semicontinuous in \(\tau\).
Its formal Legendre ray is
\[
 \ell_t(x)=
 \sup_{\tau\in\mathbf R}\{\psi_\tau(x)+t\tau\},
 \qquad t\geq0.
\tag{D.12}\label{eq:appD-test-curve-ray}
\]
\end{definition}

\begin{proposition}\label{prop:appD-test-curve-order}
If \(\tau\mapsto\psi_\tau(x)\) is proper, then
\[
 \psi_\tau(x)=\inf_{t\geq0}\{\ell_t(x)-t\tau\}.
\tag{D.13}\label{eq:appD-pointwise-inversion}
\]
Moreover, \(\psi\leq\chi\) implies
\(\mathcal L\psi\leq\mathcal L\chi\).
\end{proposition}

\begin{proof}
The identity is \Cref{thm:appD-concave-duality} at \(x\).
The order assertion follows by taking suprema in
\eqref{eq:appD-test-curve-ray}.
\end{proof}

\begin{warning}{Analytic regularization}
Formula \eqref{eq:appD-pointwise-inversion} is scalar convex analysis.
If \(\ell_t\) is replaced by an upper-semicontinuous or psh regularization
in \(x\), exact inversion is an additional maximality theorem.
\Cref{chap:rays-filtrations} supplies it only on the stated analytic
domain.
\end{warning}

\begin{example}[Single break]
\label{ex:appD-single-break}
Let \(a<b\), \(c>0\), and
\[
 g(\tau)=
 \begin{cases}
 0,&\tau\leq a,\\
 -c(\tau-a),&a\leq\tau\leq b,\\
 -\infty,&\tau>b.
 \end{cases}
\]
This is closed, concave, and nonincreasing.
Direct maximization gives a convex piecewise affine \(\mathcal Lg\) whose
slopes lie in \([a,b]\).
\end{example}

\section{Convex bodies and mixed volume}
\label{sec:appD-mixed-volume}

Assume that \(N\) carries a lattice and normalize Lebesgue measure so a
fundamental parallelepiped has volume one.
A convex body is a nonempty compact convex subset.
For \(s,t\geq0\),
\[
 sK+tL=\{sx+ty:x\in K,\ y\in L\}.
\tag{D.14}\label{eq:appD-minkowski-sum}
\]

\begin{lemma}\label{lem:appD-support-calculus}
For convex bodies \(K,L\), \(s,t\geq0\), and \(a\in N_{\mathbf R}\),
\[
\begin{aligned}
 h_{sK+tL}&=s h_K+t h_L,\\
 h_{K+a}(u)&=h_K(u)+\langle u,a\rangle,\\
 K\subseteq L&\Longleftrightarrow h_K\leq h_L.
\end{aligned}
\tag{D.15}\label{eq:appD-support-calculus}
\]
\end{lemma}

\begin{proof}
The first two formulas follow by separating suprema.
If \(K\not\subset L\), strict separation of a point of \(K\setminus L\)
from \(L\) gives \(u\) with \(h_K(u)>h_L(u)\).
\end{proof}

\begin{theorem}\label{thm:appD-minkowski-polynomial}
For convex bodies \(K_1,\ldots,K_r\subset N_{\mathbf R}\),
\[
 \operatorname{vol}\!\left(\sum_{i=1}^r t_iK_i\right)
 =
 \sum_{\substack{\alpha\in\mathbf N^r\\|\alpha|=n}}
 \frac{n!}{\alpha_1!\cdots\alpha_r!}
 V(K_1[\alpha_1],\ldots,K_r[\alpha_r])
 \prod_{i=1}^r t_i^{\alpha_i}.
\tag{D.16}\label{eq:appD-minkowski-polynomial}
\]
The coefficient \(V(K_1,\ldots,K_n)\) is symmetric, nonnegative,
translation invariant, monotone, and Hausdorff-continuous, with
\(V(K,\ldots,K)=\operatorname{vol}(K)\).
\end{theorem}

\begin{proof}
We first treat polytopes.  A fine mixed subdivision of
\(t_1K_1+\cdots+t_rK_r\), obtained from a triangulation of the Cayley
polytope, has top-dimensional cells
\[
                     t_1F_1+\cdots+t_rF_r
\]
with the direction spaces of the simplices \(F_i\) independent and
\(\sum_i\dim F_i=n\).  If \(\alpha_i=\dim F_i\), lattice-normalized
ambient volume and a linear change of variables on the direct sum of these
direction spaces give
\[
 \operatorname{vol}(t_1F_1+\cdots+t_rF_r)
       =c_F\prod_i t_i^{\alpha_i},\qquad c_F\geq0.
\]
The same subdivision works throughout the positive orthant: changing the
positive \(t_i\) only rescales the summands in every mixed cell.  Summing
over the top cells proves that the volume is a homogeneous polynomial of
degree \(n\), with nonnegative coefficients.  This mixed-subdivision
construction is the combinatorial core of
\cite[Theorem~5.1.7]{Sch14}.

For arbitrary convex bodies choose polytopes \(K_{i,j}\) converging to
\(K_i\) in Hausdorff distance.  Minkowski addition and volume are
Hausdorff-continuous, so the corresponding volume polynomials converge at
every point of \(\mathbf R_{\geq0}^r\).  The coefficients of a polynomial
of degree at most \(n\) are fixed linear combinations of its values on any
unisolvent finite grid.  They therefore converge, and their limits remain
nonnegative.  Denote the coefficient of
\(t_1^{\alpha_1}\cdots t_r^{\alpha_r}\) by
\[
       \frac{n!}{\alpha_1!\cdots\alpha_r!}
       V(K_1[\alpha_1],\ldots,K_r[\alpha_r]).
\]
This proves \eqref{eq:appD-minkowski-polynomial} and Hausdorff
continuity.

The coefficient of \(t_1\cdots t_n\) is unchanged by permuting the
bodies, which gives symmetry.  Replacing the first body by \(sK+tL\) and
comparing that coefficient on the two sides of the already proved volume
polynomial gives
\[
 V(sK+tL,K_2,\ldots,K_n)
 =sV(K,K_2,\ldots,K_n)+tV(L,K_2,\ldots,K_n).
\]
Thus mixed volume is Minkowski-multilinear.  Translating one body merely
translates every Minkowski sum, hence leaves the coefficient unchanged.
Taking every body equal to \(K\) and comparing with
\(\operatorname{vol}((\sum t_i)K)\) gives
\(V(K,\ldots,K)=\operatorname{vol}(K)\).

It remains to record monotonicity.  Polarizing the facet formula
\[
 n\operatorname{vol}(P)
    =\sum_F h_P(u_F)\operatorname{vol}_{n-1}(F)
\]
in its last \(n-1\) entries gives, first for polytopes,
\[
 nV(K,K_2,\ldots,K_n)
   =\int_{S^{n-1}}h_K\,\rd S(K_2,\ldots,K_n),
\]
where the mixed area measure is positive: on each common normal cone its
mass is the \((n-1)\)-dimensional mixed-cell coefficient already
constructed above.  Thus \(K\subseteq L\) implies \(h_K\leq h_L\) and
\(V(K,K_2,\ldots,K_n)\leq V(L,K_2,\ldots,K_n)\).
Approximating an inclusion from inside and outside by polytopes and using
the continuity just proved gives the assertion for convex bodies.  This
also completes the proof when some bodies are lower-dimensional.
\end{proof}

The theorem asserts nothing about infinite-dimensional mixed volumes.

\begin{theorem}\label{thm:appD-af}
\leavevmode\par
Assume \(n\geq2\).  For convex bodies \(K,L,K_3,\ldots,K_n\),
\[
 V(K,L,K_3,\ldots,K_n)^2
 \geq
 V(K,K,K_3,\ldots,K_n)
 V(L,L,K_3,\ldots,K_n).
\tag{D.17}\label{eq:appD-af}
\]
No nonempty-interior hypothesis is needed for the inequality.
Division by a mixed volume or an equality criterion requires the relevant
nondegeneracy hypotheses.
\end{theorem}

\begin{proof}
We give the complete reduction to one imported finite-dimensional core.
First suppose that all bodies have smooth strictly convex support
functions \(h_1,\ldots,h_n\).  Fix \(h_3,\ldots,h_n\) and, on the unit
sphere, set
\[
 Q(f,g):=\frac1{n!}\int_{S^{n-1}}
       f\,\dprimeddprime g\wedge
       \dprimeddprime h_3\wedge\cdots\wedge
       \dprimeddprime h_n\wedge\ddprime\lVert x\rVert .
\]
Stokes' formula makes \(Q\) symmetric, and the mixed-volume integral
formula gives
\[
                 Q(h_K,h_L)=V(K,L,K_3,\ldots,K_n).
\]
The imported input is exactly the following Lorentzian assertion: if
\(Q(p,p)>0\), then \(Q\) is nonpositive on
\(\{f:Q(f,p)=0\}\), with the kernel supplied by the solvability theory for
the self-adjoint Alexandrov equation.  The signature statement is
\cite[Theorem~12.3]{Ber26Convex}, and the kernel and solvability statement is
\cite[Theorem~13.3]{Ber26Convex}; the classical convex-body formulation is
\cite[Theorem~7.3.1]{Sch14}.  We use that foundational input and prove every
remaining deduction below.

Apply the Lorentzian assertion with \(p=h_L\).  When \(Q(h_L,h_L)>0\), put
\[
       f=h_K-\frac{Q(h_K,h_L)}{Q(h_L,h_L)}h_L.
\]
Then \(Q(f,h_L)=0\), hence \(Q(f,f)\leq0\).  Expanding this inequality
gives
\[
                 Q(h_K,h_L)^2
                 \geq Q(h_K,h_K)Q(h_L,h_L),
\]
which is \eqref{eq:appD-af}.  The same conclusion when the denominator
vanishes follows by replacing every body by its Minkowski sum with
\(\varepsilon B\).

Finally, every convex body is a Hausdorff limit of smooth strictly convex
bodies, for example after adding \(\varepsilon B\) and smoothing the
support function.  Mixed-volume continuity from
\Cref{thm:appD-minkowski-polynomial} lets \(\varepsilon\downarrow0\) and
proves the stated degenerate case without any division.
\end{proof}

\begin{corollary}\label{cor:appD-brunn-minkowski}
For convex bodies \(K,L\),
\[
 \operatorname{vol}(K+L)^{1/n}
 \geq
 \operatorname{vol}(K)^{1/n}
 +\operatorname{vol}(L)^{1/n}.
\tag{D.18}\label{eq:appD-brunn-minkowski}
\]
The usual homothety equality criterion requires nonempty interiors.
\end{corollary}

\begin{proof}
Put
\[
 a_j=V(K[n-j],L[j]),\qquad 0\leq j\leq n.
\]
Alexandrov--Fenchel says that \((a_j)\) is log concave.
After adding a small ball, all \(a_j\) are positive, and log concavity gives
\[
 a_j\geq a_0^{(n-j)/n}a_n^{j/n}.
\]
Insert this estimate in the Minkowski polynomial for
\(\operatorname{vol}(K+L)\) and apply the binomial theorem.
Letting the radius of the added ball tend to zero proves
\eqref{eq:appD-brunn-minkowski}.
\end{proof}

\section{Alexandrov Monge--Amp\`ere measure}
\label{sec:appD-real-ma}

Let \(\Omega\subset N_{\mathbf R}\) be open and convex and
\(u:\Omega\to\mathbf R\) convex.
For \(E\subset\Omega\), put
\(\partial u(E)=\bigcup_{x\in E}\partial u(x)\).

\begin{theorem}\label{thm:appD-alexandrov-measure}
The assignment
\[
 \MA_{\mathbf R}(u)(E)
 =\operatorname{Leb}_M(\partial u(E))
\tag{D.19}\label{eq:appD-alexandrov-measure}
\]
defines a positive Borel measure.
If \(u\in C^2(\Omega)\), then
\[
 \MA_{\mathbf R}(u)=\det(D^2u)\,\rd x.
\tag{D.20}\label{eq:appD-smooth-real-ma}
\]
If finite convex \(u_j\to u\) locally uniformly, then
\(\MA_{\mathbf R}(u_j)\) converges weakly on relatively compact subsets to
\(\MA_{\mathbf R}(u)\).
\end{theorem}

\begin{proof}
The subgradient graph is relatively closed in \(\Omega\times M\): this
follows immediately by passing to the limit in the subgradient inequality.
Hence the image
\(\partial u(E)\) of a Borel set is analytic and therefore Lebesgue
measurable.  If \(E_1,E_2\) are disjoint and
\(m\in\partial u(E_1)\cap\partial u(E_2)\), then there are distinct
\(x_i\in E_i\) with
\[
 m\in\partial u(x_i)
 \quad\Longleftrightarrow\quad
 x_i\in\partial u^*(m).
\]
Thus \(u^*\) is not differentiable at \(m\).  A finite convex function on
the relative interior of its domain is differentiable almost everywhere;
the boundary of a full-dimensional convex domain, and every
lower-dimensional domain, has \(n\)-dimensional Lebesgue measure zero.
It follows that subgradient images of disjoint Borel sets overlap only on
a null set.  Since
\(\partial u(\bigcup_jE_j)=\bigcup_j\partial u(E_j)\), the assignment in
\eqref{eq:appD-alexandrov-measure} is countably additive.

It is locally finite.  Indeed, if \(K\Subset\Omega\), choose
\(\rho>0\) with \(K+\rho B\Subset\Omega\).  For \(x\in K\),
\(m\in\partial u(x)\), and a unit vector \(e\), the subgradient inequality
at \(x\pm\rho e\) bounds \(|\langle m,e\rangle|\) by
\(2\rho^{-1}\sup_{K+\rho B}|u|\).  Hence \(\partial u(K)\) is bounded.
This proves that \(\MA_{\mathbf R}(u)\) is a positive Borel measure.

Assume next that \(u\in C^2\).  Then
\(\partial u(x)=\{\nabla u(x)\}\).  The area formula
\cite[Theorem~3.8]{EvansGariepy} gives
\[
 \int_E\det D^2u\,\rd x
 =\int_M N(\nabla u,E,m)\,\rd m.
\]
For almost every \(m\), the conjugate is differentiable and the displayed
subgradient equivalence shows that \(\nabla u(x)=m\) has at most one
solution.  The right side is therefore
\(\operatorname{Leb}_M(\nabla u(E))\), which proves
\eqref{eq:appD-smooth-real-ma}, including the possibly degenerate Hessian
case.

Finally suppose \(u_j\to u\) locally uniformly.  On each
\(K\Subset\Omega\), the preceding slope bound is uniform in \(j\).
Closedness of the subgradient inequality then gives
\[
 x_j\to x,\quad m_j\in\partial u_j(x_j),\quad m_j\to m
       \quad\Longrightarrow\quad m\in\partial u(x).
\]
Consequently, for compact \(K\), every \(\partial u_j(K)\) is eventually
contained in any prescribed neighborhood of \(\partial u(K)\), and hence
\[
 \limsup_j\MA_{\mathbf R}(u_j)(K)
       \leq\MA_{\mathbf R}(u)(K).
\]
For an open \(G\Subset\Omega\), take an \(m\in\partial u(G)\) at which
\(u^*\) is differentiable.  The function
\(x\mapsto u(x)-\langle m,x\rangle\) then has a unique minimizer
\(x\in G\).  Choose a closed ball \(D\Subset G\) centred at \(x\).
Uniqueness and compactness give a positive gap between the value at \(x\)
and the values on \(\partial D\).  Local uniform convergence therefore
puts every minimizer on \(D\) of
\(u_j-\langle m,\cdot\rangle\) in the interior of \(D\) for all large
\(j\).  Such a local minimum says \(m\in\partial u_j(G)\).  Fatou's
lemma, after discarding the null
nondifferentiability set of \(u^*\), yields
\[
 \MA_{\mathbf R}(u)(G)
       \leq\liminf_j\MA_{\mathbf R}(u_j)(G).
\]
The two Portmanteau inequalities prove weak convergence on relatively
compact subsets.  The subgradient and differentiability inputs used above
are precisely \cite[Theorems~24.5 and~25.5]{Roc70}.
\end{proof}

\begin{example}[Polyhedral atom]
\label{ex:appD-polyhedral-atom}
If
\[
 u(x)=\max_i\{\langle m_i,x\rangle+c_i\},
\]
then at a vertex \(x_0\),
\[
 \MA_{\mathbf R}(u)(\{x_0\})
 =
 \operatorname{vol}_M
 \operatorname{conv}\{m_i:
 \langle m_i,x_0\rangle+c_i=u(x_0)\}.
\tag{D.21}\label{eq:appD-polyhedral-atom}
\]
Only full-dimensional subdifferentials contribute.
\end{example}

\begin{proposition}\label{prop:appD-legendre-pushforward}
Let \(u:N_{\mathbf R}\to\mathbf R\) be closed convex and suppose
\[
 P=\overline{\partial u(N_{\mathbf R})}
\]
is compact.
Then, with \(\nabla u^*\) defined at its almost-everywhere differentiability
points in \(P\),
\[
 \MA_{\mathbf R}(u)(N_{\mathbf R})=\operatorname{vol}_M(P).
\tag{D.22}\label{eq:appD-real-ma-mass}
\]
After defining \(\nabla u^*\) arbitrarily on the Lebesgue-null subset of
\(P\) where it is not finite and differentiable, one has
\[
 (\nabla u^*)_*(\operatorname{Leb}_M|_P)
 =\MA_{\mathbf R}(u).
\tag{D.23}\label{eq:appD-legendre-pushforward}
\]
\end{proposition}

\begin{proof}
For finite convex \(u\) on all of \(N_{\mathbf R}\), the subgradient
equivalence gives
\[
 \partial u(N_{\mathbf R})=\operatorname{dom}(\partial u^*).
\]
The latter contains the relative interior of \(\operatorname{dom}(u^*)\),
and its closure is the closure of \(\operatorname{dom}(u^*)\); this closed
convex set is also \(P\).
The omitted relative boundary is Lebesgue-null, including when the domain is
lower-dimensional.  Thus the definition of Alexandrov measure gives the
first identity.  At every differentiability point of \(u^*\), the
subgradient equivalence says
\[
 \nabla u^*(m)\in E\quad\Longleftrightarrow\quad m\in\partial u(E).
\]
The nondifferentiability set is Lebesgue-null.  Applying the preceding
equivalence to Borel sets \(E\) proves the pushforward identity.
\end{proof}

\begin{warning}{The \(n!\)-test}
Our \(\MA_{\mathbf R}\) is exactly
\eqref{eq:appD-alexandrov-measure}.
For a toric line bundle,
\((L_P^n)=n!\operatorname{vol}(P)\) by
\eqref{eq:appJ-toric-degree}.
Any real-to-algebraic comparison must visibly insert \(n!\).
\end{warning}

\section{Moment measures and layer cake}
\label{sec:appD-moments}

Let \((P,\lambda)\) be a finite measure space and
\(G:P\to\mathbf R\cup\{-\infty\}\) measurable.
Its unnormalized distribution is
\[
                             \nu_G=G_*\lambda.
\tag{D.24}\label{eq:appD-pushforward-distribution}
\]
Here \(\nu_G\) is initially a measure on
\(\mathbf R\cup\{-\infty\}\).
It has mass \(\lambda(P)\) and is a probability only after division by a
positive mass.

\begin{proposition}\label{prop:appD-moment-tail}
For every nonnegative Borel
\(F:\mathbf R\cup\{-\infty\}\to[0,+\infty]\),
\[
 \int_{\mathbf R\cup\{-\infty\}}F(\tau)\,\rd \nu_G(\tau)
 =\int_PF(G(x))\,\rd \lambda(x).
\tag{D.25}\label{eq:appD-pushforward-test}
\]
If \(G\geq0\) and \(q>0\), then
\[
 \int_{\mathbf R}\tau^q\,\rd \nu_G(\tau)
 =
 q\int_0^\infty
 t^{q-1}\lambda\{G>t\}\,\rd t
\tag{D.26}\label{eq:appD-layer-cake}
\]
as an extended-real identity.
\end{proposition}

\begin{proof}
The first equality defines pushforward.
For \(a\geq0\),
\[
 a^q=q\int_0^\infty t^{q-1}\mathbf 1_{\{a>t\}}\,\rd t.
\]
Tonelli's theorem proves the second equality, including \(+\infty\).
\end{proof}

For an Okounkov concave transform, take
\(\lambda=n!\operatorname{Leb}\) to obtain the book's unnormalized
Duistermaat--Heckman measure.
Flag independence is a geometric theorem of
\Cref{chap:transcendental-okounkov-dh}, not a formal pushforward fact.

\begin{warning}{Zero mass and signed moments}
If \(\lambda(P)=0\), no probability normalization is defined.
If \(G=-\infty\) on a set of positive measure, then \(\nu_G\) has an atom at
\(-\infty\); it is not silently regarded as a Radon measure on
\(\mathbf R\), and ordinary polynomial moments are not defined.
For signed moments, treat positive and negative parts separately and never
subtract \(+\infty\) from \(+\infty\).
\end{warning}

\section{Mixed Monge--Amp\`ere and homogeneous supercurrents}
\label{sec:appD-supercurrents}

Polarization gives a symmetric positive measure
\(\MA_{\mathbf R}(u_1,\ldots,u_n)\) characterized by
\[
 \MA_{\mathbf R}\!\left(\sum_i t_i u_i\right)
 =
 \sum_{i_1,\ldots,i_n}
 t_{i_1}\cdots t_{i_n}
 \MA_{\mathbf R}(u_{i_1},\ldots,u_{i_n}),
\tag{D.27}\label{eq:appD-mixed-ma-polarization}
\]
for \(t_i\geq0\).
In the smooth case its density is the mixed discriminant of the Hessians.
Positivity and weak continuity extend it to finite convex functions.

For a convex body \(K\), put
\[
                             T_K=\dprimeddprime h_K
\tag{D.28}\label{eq:appD-support-current}
\]
in the Lagerberg convention of \Cref{chap:cld-currents}.
If \(s_t(x)=tx\) and \(E\) is radial, Euler's identity gives
\[
 s_t^*T_K=tT_K,\qquad
 \iota_E'T_K=0,\qquad
 \iota_E''T_K=0.
\tag{D.29}\label{eq:appD-strong-homogeneity}
\]
These closed positive currents are strongly homogeneous.
The phrase does not license products of arbitrary weakly positive currents.

\begin{theorem}\label{thm:appD-support-current-mixed-volume}
With lattice-normalized integration and the book's Lagerberg signs,
\[
 \frac1{n!}\,
 T_{K_1}\wedge\cdots\wedge T_{K_n}
 =
 V(K_1,\ldots,K_n)\,\delta_0.
\tag{D.30}\label{eq:appD-current-mixed-volume}
\]
The Alexandrov--Fenchel inequality is the reversed Cauchy--Schwarz
inequality for the Lorentzian form obtained after fixing
\(T_{K_3},\ldots,T_{K_n}\).
\end{theorem}

\begin{proof}
The only imported analytic input in this proof is
\cite[Theorem~4.2]{Ber26Convex}: the Bedford--Taylor product for convex
superpotentials is defined by smooth convex regularization and is continuous
under local uniform convergence.  With the book's
Berezin sign, its smooth normalization is
\[
       \frac1{n!}(\dprimeddprime u)^n
       =\det(D^2u)\,\rd x,
\]
so \Cref{thm:appD-alexandrov-measure} identifies the left side with
\(\MA_{\mathbf R}(u)\) for every finite convex \(u\).

For a support function,
\(\partial h_K(0)=K\).  If \(x\neq0\), then \(\partial h_K(x)\) is the
exposed face on which \(\langle x,\cdot\rangle\) is maximal.  If \(K\)
is full-dimensional, every such face is proper and contained in
\(\partial K\); if \(K\) is lower-dimensional, their union is contained
in \(K\).  In either case the union has Lebesgue measure zero.  Therefore
\[
               \frac1{n!}T_K^n
               =\MA_{\mathbf R}(h_K)
               =\operatorname{vol}(K)\delta_0.
\]
Now use \Cref{lem:appD-support-calculus}:
\(T_{\sum_i t_iK_i}=\sum_i t_iT_{K_i}\).  Apply the equal-body identity
to \(\sum_i t_iK_i\), expand the current power, and compare the coefficient
of \(t_1\cdots t_n\) with
\eqref{eq:appD-minkowski-polynomial}.  The coefficient on the current side
is \(T_{K_1}\wedge\cdots\wedge T_{K_n}\), while the coefficient on the
volume side is \(n!V(K_1,\ldots,K_n)\delta_0\).  This proves
\eqref{eq:appD-current-mixed-volume}.  The construction and the equivalent
mixed-volume identity are recorded in
\cite[Sections~5--7]{Ber26Convex}; that source's \(J\) and Berezin signs
have here been translated to
\Cref{chap:tropical-spaces,chap:cld-currents}.

After fixing \(K_3,\ldots,K_n\), the bilinear form represented by wedging
with \(T_{K_3}\wedge\cdots\wedge T_{K_n}\) is exactly \(n!\) times the
mixed-volume form \(Q\) used in the proof of \Cref{thm:appD-af}.
Its Lorentzian signature therefore turns \eqref{eq:appD-af} into the
reversed Cauchy--Schwarz inequality claimed in the theorem.
\end{proof}

\section{Normalization checks}
\label{sec:appD-unit-tests}

\begin{enumerate}
\item For \(f(x)=x^2/2\), one has
      \(f^*(p)=p^2/2\), \(\partial f(x)=\{x\}\), and
      \(\MA_{\mathbf R}(f)=\rd x\).
\item For \(f(x)=|x|\),
      \(f^*=\iota_{[-1,1]}\) and
      \(\MA_{\mathbf R}(f)=2\delta_0\).
\item For \(K=[0,a]\), \(L=[0,b]\),
      \(\operatorname{vol}(sK+tL)=sa+tb\).
\item Translating \(G\) by \(c\) preserves mass and changes its first moment
      by \(c\,\lambda(P)\).
\item If \(P\) has zero volume, both
      \(\MA_{\mathbf R}(h_P)\) and the unnormalized moment measure have zero
      mass, and neither can be normalized to a probability.
\end{enumerate}

The first two identities follow directly from
\eqref{eq:appD-conjugate} and \eqref{eq:appD-alexandrov-measure}; for
\(|x|\), the subdifferential at the origin is \([-1,1]\) and all other
subdifferentials are singletons.  The interval formula is ordinary
one-dimensional length.  The fourth assertion follows by testing the
pushforward against \(1\) and against the identity function.  The last
assertion is \eqref{eq:appD-real-ma-mass} together with the definition of
probability normalization.  Thus every item is a proved consequence of the
normalizations above.

\section*{Source note}

Closed convex duality and subgradients follow
\cite[Chapters~7, 12, and~23--25]{Roc70}.
Mixed volumes and Alexandrov--Fenchel follow \cite{Sch14}.
Strongly homogeneous supercurrents follow \cite{Ber26Convex}, after the sign
translation above.
Geometric realization as psh test curves or rays belongs to the cited main
chapters.
